% ============================================================
% 第1章：绪论
% ============================================================

\section{研究背景与问题定义}
电子元器件选型在硬件产品研发全生命周期中具有基础性和关键性作用。以 DC-DC 电源模块设计为例，德州仪器（Texas Instruments）、亚德诺半导体（Analog Devices）、意法半导体（STMicroelectronics）、微芯科技（Microchip Technology）等数十家制造商提供了数万种相关型号，工程师通常需要逐项核查并交叉比对其输入输出电压范围、最大输出电流、开关频率、工作温度等级、封装约束、库存水平和生命周期状态等多维参数；同时，还需查阅数据手册以验证外围元件的设计公式，并综合评估供应链稳定性 \cite{pecht2008parts}。该流程通常耗时数小时以上，且主要存在三方面瓶颈：第一，工程知识检索高度碎片化，设计规范分散于数据手册、应用笔记和工程师个人经验中；第二，供应链风险缺乏系统化的量化评估手段 \cite{chen2021supplyrisk}；第三，选型结论和输出文档格式不够规范，难以直接支撑工程评审或产品生命周期管理平台的集成应用。

针对上述问题，若能充分利用元器件管理平台的结构化物料数据与参考设计资源，并结合智能化推理技术，有望提升选型过程的自动化程度、风险识别能力和文档生成规范性。硬禾科技旗下的 eZ-PLM 是面向电子元器件全生命周期管理的开放平台（www.ezplm.cn/www.ezmpl.com），涵盖超过 110 万条在产物料数据与参考设计资源，并支持进行实时器件查询和设计案例检索，为器件检索、参考设计复用和供应链信息获取提供了数据基础。在此基础上，近年来大语言模型（LLM）与智能体（Agent）技术在工程决策辅助领域也展现出较强潜力 \cite{xi2023rise,li2023llmagent}，其自然语言理解、任务分解与交互推理能力可与领域规则引擎形成互补。因此，将 eZ-PLM 的结构化物料数据与 LLM/Agent 的智能推理能力相结合，有望将传统人工选型流程转化为智能化、可追溯的自动化决策过程。
% [插图提醒：图1-1 系统应用场景示意图——展示工程师从自然语言需求输入到获得BOM选型清单、风险评估报告、拓扑分析报告的端到端流程，左侧为用户输入框，中间为Pipeline五阶段，右侧为三类输出报告]

\section{现有方案的不足}

现有的元器件选型辅助方案大致可以分为两类。一方面，以ERP和PLM平台为代表的传统方案虽然提供了参数过滤和库存查询等基础功能，但其本质上仍依赖工程师手动设定筛选条件并逐一比对结果，既无法对器件的应用场景适配度做出智能推荐，也缺乏将供应链风险进行量化分级的能力\cite{pecht2002electronic}。另一方面，通用大语言模型（如GPT-4、DeepSeek等）虽然具备较强的自然语言理解和工程知识表达能力，但在直接回答选型问题时容易产生"幻觉"——凭空编造不存在的型号或给出缺乏数据依据的评分结论\cite{wang2023survey}。这类输出既无法溯源，也不能满足工程评审对严谨性的要求。

\section{本文贡献}

本文面向eZ-PLM平台，设计并实现了一套基于LLM与工程规则引擎协同推理的智能元器件选型与风险评估Agent系统。系统的主要技术贡献包括以下五个方面：

一是提出了五阶段线性流水线（Pipeline）架构，将"自然语言需求解析—实时API器件检索—混合智能评分—证据链构建—标准化报告生成"串联为端到端自动化流程，并辅以基于LangChain框架的推理行动（ReAct）Agent交互模式\cite{yao2022react}，同时支持批量处理与多轮对话两种使用场景；

二是设计了四级容错需求解析链，依次通过LLM语义解析、规则化类别覆盖、电压比较方向兜底和温度范围正则匹配来修正mA单位误识别、否定句式漏判、LLM拓扑幻觉等工程场景中的典型解析缺陷；

三是构建了以eZ-PLM参考设计数据为上下文约束的LLM混合评分机制，实现了"参数匹配＋供应链稳定性＋应用适配度＋设计成熟度"四维加权评分，同时支持无参考设计时自动降级为纯规则双维评分，保证了评分结果的可降级性与可比性；

四是建立了规则引擎事实验证与LLM工程叙述分离的双层风险评估架构——规则侧负责检测供应商集中度、零库存、停产等客观风险，LLM侧仅对已验证的风险条目进行工程背景阐述，确保了风险报告的严谨性与数据可溯源性；

五是定义了PartIR、ScoreBreakdown、TopologyIR、RiskIR等标准化中间表示（IR）数据模型，三类工程输出文档（物料清单（BOM）选型清单、供应链风险评估报告、电路拓扑分析报告）的格式可直接用于工程评审及eZ-PLM平台对接。
